\documentclass[10pt,a4paper]{amsart}
\usepackage[margin=19mm]{geometry}
\usepackage{amsmath,amssymb,mathtools}
\usepackage[scr=rsfs,cal=euler]{mathalfa}
\usepackage{xfrac}
\usepackage[unicode,hidelinks,bookmarks=false]{hyperref}
\newcommand{\ie}{i.e.\ }
\newcommand{\cf}{cf.\ }
\newcommand{\resp}{respectively\ }
\newcommand{\Subsection}{\subsection}
\providecommand{\nobreakdash}{\nobreak\hbox{-}}
\setlength{\emergencystretch}{2em}
\multlinegap0pt
\usepackage{amssymb}
\usepackage{xcolor}
\newtheorem{proposition}{Proposition}
\newcommand{\CSL}{\mathsf{CSL}}
\newcommand{\NN}{\mathbb{N}}
\newcommand{\RCA}{\mathsf{RCA}}
\newcommand{\TT}{\mathsf{TT}}
\newcommand{\USub}{\mathsf{USub}}
\newcommand{\bstr}{2^{<\NN}}
\title{$\Pi^0_4$ conservation of a Carlson-Simpson lemma\\ for 1-variable words}
\author{Quentin Le Hou\'erou, Ludovic Patey}
\date{Source: arXiv:2607.28116v1 (CC BY 4.0)}
\begin{document}
\maketitle

\noindent\emph{Source and licence.} $\Pi^0_4$ conservation of a Carlson-Simpson lemma\\ for 1-variable words. Version arXiv:2607.28116v1 (CC BY 4.0), distributed by arXiv under the Creative Commons Attribution 4.0 International licence, http://creativecommons.org/licenses/by/4.0/, as recorded in the arXiv licence metadata for this exact version and shown on the abstract page https://arxiv.org/abs/2607.28116v1. Full licence notice: \texttt{licenses/arxiv-2607.28116-CC-BY-4.0.txt}.\par\smallskip

\noindent\textit{Editorial scope.} Counted unit: the article's proposition proposition (source0.tex lines 1087--1089). The article's complete original proof of it is delivered with it (source0.tex lines 1090--1105, one \texttt{proof} environment of 16 lines). No mathematics is written, repaired or re-derived: the statement and the proof are the article's own lines, in the article's own notation, and no retained line is substituted or re-flowed.  No further source-local context is needed: every cross-reference of every retained line resolves inside the two delivered spans.  Nothing was omitted from any retained span, so the package carries no omission record.  No retained line contains a citation, so the wrapper carries no bibliography and no bibliography entry.  No retained line contains a figure environment, an image-inclusion command or drawing code, so no image is delivered and the package declares no asset dependency.  Editorial formatting only, in addition: an AMS \texttt{amsart} preamble that restates the article's own declarations and macros used by the retained lines, namely the environments \texttt{proposition} on separate counters, together with the standard packages those lines need.  The retained environments print in this document as Proposition 1 in order of appearance, whereas the article prints the same items with the numbering of its own full document.  The complete native source of the article as distributed in the arXiv e-print for this exact version is bundled unchanged as \texttt{sources/206387/source0.tex}.
\begin{proposition}
$\RCA_0 \vdash \forall k (\CSL^1_k \to \TT^2_k)$.
\end{proposition}

\begin{proof}
Let $f : \bstr \to k$ be an instance of $\TT^2_k$ for some~$k \geq 1$.
Let $A = \{0,1\}$ and $h : A^{<\omega, 1} \to k$ be defined by
$$h(v) = f(v[\epsilon], v[0])$$
By $\CSL^1_k$, there is an infinite $\omega$-variable word~$W$ over~$A$ such that
$\USub^{1,\star}_A(W)$ is $h$-monochromatic for some color~$c$. Let 
$$T = \{ \sigma \in \bstr : |\sigma| \text{ is even } \wedge \forall i < \frac{|\sigma|}{2}, \sigma(2i) = 0 \}$$
Note that $T \cong \bstr$.
Finally, let $S = \{ W[\sigma] : \sigma \in T \}$.
Since $\sigma \mapsto W[\sigma]$ is order-preserving, $S \cong \bstr$.
We claim that $[S]^2$ is $f$-monochromatic for color~$c$.
Fix some $\hat \sigma, \hat \tau \in S$ such that $\hat \sigma \prec \hat \tau$.
Let $\sigma, \tau \in T$ be such that $\hat \sigma = W[\sigma]$ and $\hat \tau = W[\tau]$. In particular, $\sigma \prec \tau$. By definition of~$T$, there is some~$\tau' \in \bstr$ such that $\tau = \sigma \cdot 0 \cdot \tau'$.
Let $u = \sigma \cdot \star \cdot \tau'$. By choice of~$W$, $h(W[u]) = c$.
By definition of~$h$, $$f(\hat \sigma, \hat \tau) = f(W[\sigma], W[\tau]) = f(W[u][\epsilon], W[u][0]) = h(W[u]) = c$$
\end{proof}

\end{document}
